\scanpage{15}
montrer que la donnée de descente envisagée est effective.

Or on \deleted{se ramène facilement au cas où $S$ est local, pt fermé $s$, $U=S-\{s\}$. Par factorisation (\unsure{à cause de} \deleted{l'hypothèse} \inserted{le fait} que $S'$ normal) on est ramené à un morphisme $S'\longrightarrow S$ qui est un épimorphisme effectif.
\inserted{$(\operatorname{Spec}R'\unsure{-|\longrightarrow}\operatorname{Spec}R)$}
\deleted{\unsure{avec}} \deleted{$R'/R$ de \unsure{deg.} $1$ sur $R$ (?).}
\unsure{Alors}, \uline{\unsure{à} l'effectivité des descentes en \unsure{fcts} de $\unsure{F_{ab}}$} \unsure{on} est ramené à montrer qu'il \deleted{\unsure{le}} \deleted{faut \unsure{alors}} \deleted{\unsure{qu'}} existe sur $A'$ une polarisation stable par \unsure{donnée de descente}}
-- un th. de descente qui nous assure cela, cf plus bas.
\margin{\figure{\begin{tikzpicture}\node[draw,ellipse,text width=2.8cm,align=center,inner sep=5pt] {\deleted{\unsure{quant à} l'effectivité \illegible{2} \unsure{à éclaircir}.}};\end{tikzpicture}}}
\mspage{14}
\figure{\begin{tikzpicture}\draw (0,0)--(13,.4);\node[anchor=north west] at (0,-.02) {?!};\draw (0,-.8)--(1.5,-.06);\end{tikzpicture}}

 e) Cas général. \deleted{Le cas} On a déjà vu qu'on peut supposer $S$ réduit. Le cas $S$ \uline{local} relève de d). Cela permet donc, pour $\forall s\in S$, de trouver un schéma abélien \inserted{$A(s)$ sur} $\operatorname{Spec}\underline O_{S,s}\inserted{=S(s)}$, qui soit une solution au Pb induit sur $S(s)$ (avec $U(s)=U\times_S S(s)$, $T(s)=T\cap S(s)$ \unsure{etc} $\ldots$). Par unicité, les $A(s)$ forment un système local de schémas abéliens, i.e. pour $t$ gén. de $s$, on a un isom can $\unsure{A(t)}\simeq A(s)\times_{S(s)}S(t)$, on a transitivité. \deleted{Supposons \unsure{qu'on puisse}} Or on a un foncteur \uline{pl. fid} \unsure{Nat} de la catégorie des schémas de t.f. sur $S$ dans la cat. des systèmes locaux de schémas de t.f. sur les $S(s)$. Supposons qu'on ait prouvé que $\{A(s)\}_s$ est dans l'image essentielle. Alors on a donc une solution \deleted{\unsure{abél.}} $A$ sur $S$, \unsure{néce-}
\unsure{abélien} $\ldots$ (par la pleine fidélité précédente). On a

\margin{\uline{Lemme} (\unsure{indépendant} de l'hypo. faite dans la partie $2^{\text{ième}}$ du th): pour que $(A_U,M,\varphi)$ \unsure{appartienne} à l'image essentielle, il f. et s. que pour $\forall$ pt \deleted{fermé} de $T$, il y ait une solution sur $\operatorname{Spec}\unsure{O_{S,t}}$. Mais \unsure{$=$|si} on se borne aux $t$ \uline{fermés}, on \unsure{aura à} choisir avec $t$ maximal dans \unsure{les} \unsure{tels points} \unsure{qui} \illegible{2} \unsure{qu'il} \unsure{y ait une solution} \unsure{dans} $\operatorname{Spec}\unsure{O_{S,t'}}$ \unsure{$-U'$}).}
